\scanpage{55}

\deleted{notons que pour dans le cas qui nous occupe,}

Mais revenons sur la définition de $Z$ tantôt, en définissant un \inserted{autre} groupe, plus petit encore a priori

\[
 Z_0=\left\{\lambda\in\unsure{\mathfrak S|\mathscr S}_S\ \middle|\
 \begin{aligned}
 &\lambda\text{ centralise }G_0=\operatorname{Aut}(S,\Sigma_0),\\
 &\exists f_\lambda\in\mathfrak S_S\text{ tel qu'on ait, }
 \deleted{\unsure{\forall\{s,t\}\in\mathscr P_2(S)}}\inserted{\forall(s,t)\in S\times S,}\\
 &\{s,t\}\in A_0\Longleftrightarrow
 \{\underbrace{f_\lambda(\deleted{\unsure{x}}s)}_{s'},\underbrace{\lambda^{-1}f_\lambda(\deleted{\unsure{y}}t)}_{t'}\}\in A_0
 \end{aligned}\right.
 \tag*{(138)}
\]

en d'autres termes, $\lambda\in Z_0$ \unsure{sss|ssi} \deleted{pour} \inserted{$\lambda$ centralise $G_0=\operatorname{Aut}(S,\Sigma_0)$, et si pour} $s',t'\in S$, \deleted{la relation} on a

\[
 \bigl(\{s',\lambda^{-1}t'\}\in A_0\Longleftrightarrow
 \{\lambda^{-1}s',t'\}\in A_0\bigr)\Longrightarrow(s'\ne t')
 \tag*{(139)}
\]

de sorte qu'on peut \deleted{désigner par $A$} \inserted{introduire}

\[
 A_\lambda=\left\{\{s',t'\}\in\unsure{\mathscr P|\mathcal P|\mathfrak P}_2(S)\ \middle|\
 \begin{aligned}
 &\{s',\lambda^{-1}t'\}\in A_0\quad\text{i.e.}\\
 &\exists\unsure{(s,t)|\{s,t\}}\in S^2,\ \text{avec }\{s,t\}\in A_0,\\
 &\{s',\deleted{\unsure{s}}t'\}=\{s,\lambda.t\}
 \end{aligned}\right\}
\]

et si $A_\lambda$ est l'ens des arêtes \deleted{pour} une structure triangulée sur $S$ \deleted{\uline{équivalente}} \inserted{isomorphe} à la structure $\Sigma_0$ définie par $A_0\subset\mathscr P_2(S)$, i.e. s'il existe $f_\lambda\in\mathfrak S_S$ telle qu'on ait

\[
 f_\lambda(A_0)=A_\lambda
 \tag*{(140)}
\]

i.e. justement tel que

\[
 \{s,t\}\in A_0\Longleftrightarrow\{f_\lambda(s),\lambda^{-1}f_\lambda(t)\}\in A_0.
 \tag*{(141)}
\]

% archive sheet number: 54
